% Completeness20261024: derive the four-vector relations in a fixed metric convention.
下面讨论静质量 $m>0$ 恒定的质点，采用 $+---$ 度规。质点自身时钟走过的固有时满足
\[
 c^2\dd\tau^2=c^2\dd t^2-\dd\vect r^2,
 \qquad \dd\tau=\frac{\dd t}{\gamma},\qquad
 \gamma=(1-v^2/c^2)^{-1/2}.
\]
四位移除以不变的固有时，得到四速度
\[
 U^\mu=\frac{\dd x^\mu}{\dd\tau}=\gamma(c,\vect v),
 \qquad U^\mu U_\mu=\gamma^2(c^2-v^2)=c^2.
\]
再乘静质量得四动量 $p^\mu=mU^\mu=(E/c,\vect p)$，所以
\[
 E=\gamma mc^2,\quad \vect p=\gamma m\vect v,\quad
 E^2-p^2c^2=m^2c^4,\quad \vect v=\frac{c^2\vect p}{E}.
\]
能量是四动量的时间分量乘 $c$，不是另外拼接的标量。
在低速极限 $E=mc^2+mv^2/2+O(v^4/c^2)$，扣去静能后回到经典动能。
对无静质量粒子应直接用 $E=pc$ 的四动量关系，不能使用质点固有时定义四速度。

\paragraph{功率关系和四力的时间分量}
对能量--动量关系微分，并利用 $m$ 不变，得到
\[
 E\dd E=c^2\vect p\cdot\dd\vect p,
 \qquad \dd E=\vect v\cdot\dd\vect p.
\]
定义三维力 $\vect F=\dd\vect p/\dd t$，便有
$\dd E/\dd t=\vect F\cdot\vect v$。
由于 $\dd/\dd\tau=\gamma\dd/\dd t$，四力为
\[
 K^\mu=\frac{\dd p^\mu}{\dd\tau}
 =\gamma\left(\frac{\vect F\cdot\vect v}{c},\vect F\right).
\]
因而 $\vect F\cdot\dd\vect r=\dd E$ 仍成立；相对论改变的是动量与速度的关系，
不能在任意方向直接写成 $\vect F=m\vect a$。

\paragraph{四加速度为什么与四速度正交}
令 $A^\mu=\dd U^\mu/\dd\tau$。对 $U^\mu U_\mu=c^2$ 求导，立刻得
\[
 2U_\mu A^\mu=0.
\]
这里是Minkowski内积为零，不是三维速度与三维加速度垂直。
因 $m$ 恒定，$K^\mu=mA^\mu$，故也有 $U_\mu K^\mu=0$。
若物体吸收能量而使静质量改变，则
$K^\mu=mA^\mu+(\dd m/\dd\tau)U^\mu$，不能照搬后一正交关系。

写 $\vect a=\dd\vect v/\dd t$，由
$\dot\gamma=\gamma^3\vect v\cdot\vect a/c^2$ 可逐分量求得
\[
 A^0=\gamma^4\frac{\vect v\cdot\vect a}{c},\qquad
 \vect A=\gamma^2\vect a+
 \gamma^4\frac{\vect v\cdot\vect a}{c^2}\vect v.
\]
取图~\ref{comp:A6:proper-acceleration}的 $x$ 轴沿瞬时速度，便有
$A^0=\gamma^4va_\parallel/c$、$A^x=\gamma^4a_\parallel$、
$A^y=\gamma^2a_\perp$。
向速度同为 $v$ 的瞬时本征惯性系作Lorentz变换：
\[
 \begin{aligned}
 A_0^0&=\gamma(A^0-\beta A^x)=0,\\
 A_0^x&=\gamma(A^x-\beta A^0)=\gamma^3a_\parallel,\\
 A_0^y&=A^y=\gamma^2a_\perp.
 \end{aligned}
\]
在这一瞬间本征系中的速度为零，四加速度的空间部分就是该系测得的加速度，故
\[
 \boxed{a_{\parallel,0}=\gamma^3a_\parallel,\qquad
 a_{\perp,0}=\gamma^2a_\perp},\qquad
 -A^\mu A_\mu=a_0^2=\gamma^6a_\parallel^2+\gamma^4a_\perp^2.
\]
\begin{center}\begin{minipage}{.94\linewidth}\centering
\begin{tikzpicture}[>=Latex,font=\small,scale=.9]
 \begin{scope}
  \draw[->] (0,0)--(3.3,0) node[right]{$x\parallel\vect v$};
  \draw[->] (0,0)--(0,2.3) node[above]{$y$};
  \draw[->,very thick,DeepBlue] (0,0)--(2,1.4) node[above right]{$\vect a$};
  \draw[dashed] (2,0)--(2,1.4)--(0,1.4);
  \node[below] at (1,0) {$a_\parallel$};\node[left] at (0,.8) {$a_\perp$};
  \node[below] at (1.3,-.7) {地面系：$\vect v\ne0$};
 \end{scope}
 \begin{scope}[xshift=7.2cm]
  \draw[->] (0,0)--(3,0) node[right]{$x_0$};
  \draw[->] (0,0)--(0,2.3) node[above]{$y_0$};
  \draw[->,very thick,DeepBlue] (0,0)--(2.25,1.2) node[above right]{$\vect a_0$};
  \draw[dashed] (2.25,0)--(2.25,1.2)--(0,1.2);
  \node[below] at (1.25,0) {$\gamma^3a_\parallel$};
  \node[left] at (0,.7) {$\gamma^2a_\perp$};
  \node[below] at (1.3,-.7) {瞬时本征系：$\vect v_0=0$};
 \end{scope}
\end{tikzpicture}
\figcap{同一事件的加速度分解。两幅图各用示意比例；平行和垂直分量的变换因子不同，方向一般也改变。}
\label{comp:A6:proper-acceleration}
\end{minipage}\end{center}
图中用相同的平行、垂直方向识别分量；它不是在同一惯性系中把两根加速度箭头相加。
例如地面系中若 $a_\parallel=a_\perp$，本征系的比值就为
$a_{\perp,0}/a_{\parallel,0}=1/\gamma$，本征加速度更靠近平行方向。

\paragraph{用快度把共线合成化为加法}
定义 $\vartheta$ 使 $\beta=\tanh\vartheta$，则
$\gamma=\cosh\vartheta$、$\gamma\beta=\sinh\vartheta$，
$U^\mu=c(\cosh\vartheta,\sinh\vartheta,0,0)$。
对同向共线速度的合成，
\[
 \beta_{12}=\frac{\beta_1+\beta_2}{1+\beta_1\beta_2}
 =\tanh(\vartheta_1+\vartheta_2),
\]
所以快度相加；反向相对运动按有向快度相减。其依据是双曲正切的加法恒等式，
而不是把接近光速的普通速度直接相加。
